\chapter{Segunda Parte: Metodologia e Medições Experimentais}

\subsection{Instrumentação e Equipamentos Utilizados}
A confiabilidade das medições experimentais desta prática repousa fundamentalmente na precisão e calibração dos instrumentos de bancada empregados. Todo o procedimento laboratorial de excitação, leitura e diagnóstico dos parâmetros elétricos foi garantido por meio do seguinte \textit{setup} de alto desempenho:

\begin{itemize}
	\item \textbf{Multímetro Digital (True RMS): Keysight U1253B.} Utilizado primariamente para a aferição a frio dos componentes passivos (medição ôhmica e capacitiva) com alta acurácia, atenuando as incertezas paramétricas na posterior análise matemática dos dados.
	\item \textbf{Fonte de Tensão Contínua: Keysight E3631A.} Responsável pela polarização do circuito emissor comum, fornecendo a tensão contínua assimétrica de $V_{CC} = 12\text{ V}$, garantindo baixos níveis de \textit{ripple} e ruído de linha acoplado.
	\item \textbf{Osciloscópio Digital (DSO): Keysight DSO-X 2012A.} Operado no modo de múltiplos canais simultâneos para o monitoramento síncrono do sinal de entrada (Canal 1) e da resposta do amplificador (Canal 2), possibilitando a leitura gráfica de amplitudes, análise de defasagem e a constatação de distorções na saída (como o possível ceifamento devido à saturação do componente).
	\item \textbf{Gerador de Funções do osciloscópio:} Empregado para injetar o sinal senoidal de excitação ($v_{sig}$) na entrada do circuito, configurado rigorosamente para a frequência de $5\text{ kHz}$ com amplitude ajustada para refletir $30\text{ mV}$ de pico a pico no primeiro nó de teste.
\end{itemize}

\subsection{Aferição a Frio dos Componentes Passivos}
Para garantir a fidelidade da estimativa numérica (via MMEQ) na fase pós-laboratorial, os resistores e capacitores comerciais empregados na montagem da matriz de contatos foram previamente medidos a frio. O registro exato permite mitigar os erros associados às tolerâncias de fábrica dos componentes (tipicamente na faixa de $5\%$ a $10\%$). Os valores efetivos encontram-se sumarizados na Tabela \ref{tab:afericao_frio}.

\begin{table}[htbp]
	\centering
	\caption{Valores nominais e medidos a frio dos componentes passivos do circuito.}
	\label{tab:afericao_frio}
	\begin{tabular}{lcc}
		\toprule
		\textbf{Componente} & \textbf{Valor Nominal} & \textbf{Valor Medido} \\
		\midrule
		Resistor de Coletor ($R_C$) & $4,7\text{ k}\Omega$ & $4,68\text{ k}\Omega$ \\
		Resistor de Emissor ($R_E$) & $470\ \Omega$ & $465,01\ \Omega$ \\
		Resistor de Realimentação ($R_f$) & $150\text{ k}\Omega$ & $149,3\text{ k}\Omega$ \\
		Resistor de Medição 1 ($R_{m1}$) & $220\ \Omega$ & $220,21\ \Omega$ \\
		Resistor de Medição 2 ($R_{m2}$) & $680\ \Omega$ & $684,01\ \Omega$ \\
		Potenciômetro (Carga máxima) & $50\text{ k}\Omega$ & $53,6\text{ k}\Omega$ \\
		\midrule
		Capacitor de Acoplamento 1 ($C_{C1}$) & $10\ \mu\text{F}$ & $9,44\ \mu\text{F}$ \\
		Capacitor de Acoplamento 2 ($C_{C2}$) & $10\ \mu\text{F}$ & $12,30\ \mu\text{F}$ \\
		Capacitor de Desvio ($C_E$) & $10\ \mu\text{F}$ & $12,43\ \mu\text{F}$ \\
		\bottomrule
	\end{tabular}
\end{table}